\documentclass[10pt,a4paper]{amsart}
\usepackage[margin=19mm]{geometry}
\usepackage{amsmath,amssymb,mathtools}
\usepackage[scr=rsfs,cal=euler]{mathalfa}
\usepackage{xfrac}
\usepackage[unicode,hidelinks,bookmarks=false]{hyperref}
\newcommand{\ie}{i.e.\ }
\newcommand{\cf}{cf.\ }
\newcommand{\resp}{respectively\ }
\newcommand{\Subsection}{\subsection}
\providecommand{\nobreakdash}{\nobreak\hbox{-}}
\setlength{\emergencystretch}{2em}
\multlinegap0pt
\usepackage{amssymb}
\usepackage{xcolor}
\newtheorem{theorem}{Theorem}
\newtheorem{corollary}{Corollary}
\title{The discrete second moment of mixed derivatives of the Riemann zeta function}
\author{Benjamin Durkan, Christopher Hughes,, Andrew Pearce-Crump}
\date{Source: arXiv:2601.06292v1 (CC BY 4.0)}
\begin{document}
\maketitle

\noindent\emph{Source and licence.} The discrete second moment of mixed derivatives of the Riemann zeta function. Version arXiv:2601.06292v1 (CC BY 4.0), distributed by arXiv under the Creative Commons Attribution 4.0 International licence, http://creativecommons.org/licenses/by/4.0/, as recorded in the arXiv licence metadata for this exact version and shown on the abstract page https://arxiv.org/abs/2601.06292v1. Full licence notice: \texttt{licenses/arxiv-2601.06292-CC-BY-4.0.txt}.\par\smallskip

\noindent\textit{Editorial scope.} Counted unit: the article's theorem mainthm (source0.tex lines 120--157). The article's complete original proof of it is delivered with it (source0.tex lines 733--807, one \texttt{proof} environment of 75 lines, which the article places away from the statement). No mathematics is written, repaired or re-derived: the statement and the proof are the article's own lines, in the article's own notation, and no retained line is substituted or re-flowed.  Retained uncounted as the source-local context the proof needs: lines 95--97; lines 171--173.  Nothing was omitted from any retained span, so the package carries no omission record.  No retained line contains a citation, so the wrapper carries no bibliography and no bibliography entry.  No retained line contains a figure environment, an image-inclusion command or drawing code, so no image is delivered and the package declares no asset dependency.  Editorial formatting only, in addition: an AMS \texttt{amsart} preamble that restates the article's own declarations and macros used by the retained lines, namely the environments \texttt{theorem}, \texttt{corollary} on separate counters, together with the standard packages those lines need.  The retained environments print in this document as Theorem 1, Corollary 1 in order of appearance, whereas the article prints the same items with the numbering of its own full document.  The complete native source of the article as distributed in the arXiv e-print for this exact version is bundled unchanged as \texttt{sources/192398/source0.tex}.
    \begin{equation}\label{milinovich_theorem}
        \sum_{0<\gamma\le T}\left|\zeta'\left(\frac{1}{2}+i\gamma\right)\right|^2=\frac{T}{2\pi} P_4\left(\log\frac{T}{2\pi}\right)+O\left(T^{\frac{1}{2}+\varepsilon}\right)
    \end{equation}

\begin{corollary} \label{cor:micah}
    Assume the Riemann Hypothesis. In the case $\mu=\nu=1$, \eqref{mainthmeq} recovers the polynomial $P_4(x)$ in \eqref{milinovich_theorem}.
\end{corollary}

\begin{theorem}\label{mainthm}
    For positive integers $\mu,\nu$, we have 
    \begin{equation}\label{mainthmeq}
        \sum_{0<\gamma\le T}\zeta^{(\mu)}(\rho)\zeta^{(\nu)}(1-\rho)=\frac{T}{2\pi}\mathcal{P}_{\mu,\nu}\left(\log\frac{T}{2\pi}\right)+O\left(Te^{-C\sqrt{\log T}}\right)
    \end{equation}
    as $T\to\infty$, where $C$ is a positive constant and where $\mathcal{P}_{\mu,\nu}(x)$ is the polynomial of degree $\mu+\nu+2$ given by
    \begin{multline}\label{eq:poly}
        \mathcal{P}_{\mu,\nu}(x) = \sum_{m=0}^{\mu+\nu+2} \sum_{k=0}^{\nu}  (-1)^{\nu}  \binom{\nu}{k} C_1^{(\mu,\nu)}(m,k)x^m\\
        +\sum_{m=0}^{\mu+\nu+2} \sum_{k=0}^{\mu}  (-1)^{\mu}  \binom{\mu}{k} \left(C_1^{(\nu,\mu)}(m,k)+C_2^{(\mu,\nu)}(m,k)\right)x^m,
    \end{multline}
    where
\begin{multline*}
C_1^{(\mu,\nu)}(m,k) = \\
\begin{cases} 
\displaystyle \frac{(-1)^m(\nu-k)}{m!} \sum_{j=0}^{\mu+\nu+1-m}  (-1)^{\mu+\nu-j} \frac{ (\mu+\nu+1-j)! }{(\mu+k+2-j)!} c^{(\mu,k)}_{j} +\\
\displaystyle \qquad + \frac{c^{(\mu,k)}_{\mu+\nu+2-m}}{(k+m-\nu)!} &  m \geq \nu-k \\[3ex]
\displaystyle \frac{(-1)^m(\nu-k)}{m!} \sum_{j=0}^{\mu+k+2}  (-1)^{\mu+\nu-j} \frac{ (\mu+\nu+1-j)! }{(\mu+k+2-j)!} c^{(\mu,k)}_{j}  & m \leq \nu -k-1
\end{cases}
\end{multline*}
where $c^{(\mu,k)}_{j}$ are the Laurent series coefficients around $s=1$ of \begin{equation}
    \frac{\zeta'}{\zeta}(s)\zeta^{(\mu)}(s)\zeta^{(k)}(s)\frac{1}{s}=\sum_{j=0}^{\infty}c^{(\mu,k)}_{j}(s-1)^{-\mu-k-3+j} ,
\end{equation}
    and  where
\begin{multline*}
    C_2^{(\mu,\nu)}(m,k) = \\
\begin{cases} 
\displaystyle  \frac{(-1)^m(\mu-k+1)}{m!} \sum_{j=0}^{\mu+\nu+1-m} (-1)^{\mu+\nu-j}  \frac{(\mu+\nu+1-j)!}{(\nu+k+1-j)!} d_{j}^{(\nu,k)} +  \\
\displaystyle \qquad  +\frac{d_{\mu+\nu+2-m}^{(\nu,k)}}{(k+m-\mu-1)!} & m \geq \mu-k+1 \\[3ex]
\displaystyle  \frac{(-1)^m(\mu-k+1)}{m!} \sum_{j=0}^{\nu+k+1} (-1)^{\mu+\nu-j}  \frac{(\mu+\nu+1-j)!}{(\nu+k+1-j)!} d_{j}^{(\nu,k)}  & m \leq \mu-k
\end{cases}
\end{multline*}
where $d_{j}^{(\nu,k)}$ are the Laurent series coefficients around $s=1$ of 
    \begin{equation*}
        \zeta^{(\nu)}(s)\zeta^{(k)}(s)\frac{1}{s}=\sum_{j=0}^{\infty}d_{j}^{(\nu,k)}(s-1)^{-\nu-k-2+j}.
    \end{equation*}

If one assumes the Riemann Hypothesis, the error term may be replaced with $O\left(T^{\frac{1}{2}+\varepsilon}\right)$ for arbitrary $\varepsilon>0$.
\end{theorem}

\begin{proof}[Proof of Corollary \ref{cor:micah}]

We show that our theorem with $\mu=\nu=1$ recovers precisely Milinovich's full asymptotic expansion for the first derivative, stated in \eqref{milinovich_theorem}.

Theorem~\ref{mainthm} states that, under the Riemann Hypothesis,
\begin{multline*}
\sum_{0<\gamma\le T}\left|\zeta'\left(\frac{1}{2}+i\gamma\right)\right|^2 \\
= \frac{T}{2\pi} \sum_{m=0}^4 (-2C_1^{(1,1)}(m,0) -C_2^{(1,1)}(m,0) -  2C_1^{(1,1)}(m,1) - C_2^{(1,1)}(m,1)) \left(\log\frac{T}{2\pi}\right)^m \\
+ O\left(T^{\frac{1}{2}+\varepsilon}\right).
\end{multline*}

We now calculate the various terms in this expansion. We have
\[
C_1^{(1,1)}(m,0) = 
\begin{cases}
     c^{(1,0)}_0-c^{(1,0)}_1+c^{(1,0)}_2-c^{(1,0)}_3 & \text{ if } m=0 \\[1.5ex]
    -c^{(1,0)}_0+c^{(1,0)}_1-c^{(1,0)}_2+c^{(1,0)}_3 & \text{ if } m=1 \\[1.5ex]
   \frac{1}{2} c^{(1,0)}_0- \frac{1}{2}  c^{(1,0)}_1 +c^{(1,0)}_2 & \text{ if } m=2 \\[1.5ex]
    -\frac{1}{6} c^{(1,0)}_0 +\frac{1}{2} c^{(1,0)}_1 & \text{ if } m=3\\[1.5ex]
    \frac{1}{6} c^{(1,0)}_0 & \text{ if } m=4
\end{cases}
\]
where $c^{(1,0)}_j$ are the coefficients of
\[
\frac{\zeta'(s)^2}{s} = \frac{c^{(1,0)}_0}{(s-1)^4} + \frac{c^{(1,0)}_1}{(s-1)^3} + \frac{c^{(1,0)}_2}{(s-1)^2} + \frac{c^{(1,0)}_3}{(s-1)} + \dots
\]
and where
\[
C_2^{(1,1)}(m,0) = 
\begin{cases}
6 d^{(1,0)}_0 -4 d^{(1,0)}_1 + 2d^{(1,0)}_2 & \text{ if } m=0 \\[1.5ex]
-6 d^{(1,0)}_0 + 4d^{(1,0)}_1 - 2d^{(1,0)}_2 & \text{ if } m=1 \\[1.5ex]
3 d^{(1,0)}_0 -2d^{(1,0)}_1 +d^{(1,0)}_2& \text{ if } m=2 \\[1.5ex]
-d^{(1,0)}_0 + d^{(1,0)}_1& \text{ if } m=3\\[1.5ex]
\frac{1}{2} d^{(1,0)}_0 & \text{ if } m=4
\end{cases}
\]
where $d^{(1,0)}_j$ are the coefficients of
\[
\frac{\zeta'(s) \zeta(s)}{s} = \frac{d^{(1,0)}_0}{(s-1)^3} + \frac{d^{(1,0)}_1}{(s-1)^2} + \frac{d^{(1,0)}_2}{(s-1)} + \dots
\]
and where
\[
C_1^{(1,1)}(m,1) = 
\begin{cases}
c^{(1,1)}_4 & \text{ if } m=0 \\[1.5ex]
c^{(1,1)}_3 & \text{ if } m=1 \\[1.5ex]
\frac{1}{2} c^{(1,1)}_2 & \text{ if } m=2 \\[1.5ex]
\frac{1}{6} c^{(1,1)}_1 & \text{ if } m=3\\[1.5ex]
\frac{1}{24} c^{(1,1)}_0 & \text{ if } m=4
\end{cases}
\]
where $c^{(1,1)}_j$ are the coefficients of
\[
\frac{\zeta'(s)^3}{\zeta(s)} \frac{1}{ s} = \frac{c^{(1,1)}_0}{(s-1)^5} + \frac{c^{(1,1)}_1}{(s-1)^4} + \frac{c^{(1,1)}_2}{(s-1)^3} + \frac{c^{(1,1)}_3}{(s-1)^2} + \frac{c^{(1,1)}_4}{(s-1)} + \dots
\]
and where
\[
C_2^{(1,1)}(m,1) = 
\begin{cases}
d^{(1,1)}_0 -d^{(1,1)}_1 +d^{(1,1)}_2 -d^{(1,1)}_3  & \text{ if } m=0 \\[1.5ex]
-d^{(1,1)}_0 +d^{(1,1)}_1 -d^{(1,1)}_2 +d^{(1,1)}_3  & \text{ if } m=1 \\[1.5ex]
\frac{1}{2} d^{(1,1)}_0 -\frac{1}{2} d^{(1,1)}_1 +d^{(1,1)}_2 & \text{ if } m=2 \\[1.5ex]
 -\frac{1}{6} d^{(1,1)}_0 +\frac{1}{2} d^{(1,1)}_1 & \text{ if } m=3 \\[1.5ex]
\frac{1}{6} d^{(1,1)}_0 & \text{ if } m=4 \\[1.5ex]
\end{cases}
\]
where $d^{(1,1)}_j$ are the coefficients of
\[
\frac{\zeta'(s)^2}{s} = \frac{d^{(1,1)}_0}{(s-1)^4} + \frac{d^{(1,1)}_1}{(s-1)^3} + \frac{d^{(1,1)}_2}{(s-1)^2} + \frac{d^{(1,1)}_3}{(s-1)} + \dots
\]
which we note are the same as $c^{(1,0)}_j$ in this special case.

Evaluating the Laurent coefficients $c^{(1,k)}_j$ and $d^{(1,k)}_j$ and inserting them into $C_1(m,k)$ and $C_2(m,k)$ yields \eqref{milinovich_theorem}.
\end{proof}

\end{document}
